\documentclass[12pt]{article}
\usepackage[a4paper,margin=1in]{geometry}
\usepackage{graphicx}
\usepackage{amsmath}
\usepackage{parskip}

\begin{document}

\begin{center}
    {\LARGE \textbf{A Review on NRAS and Kinase Inhibitors: Therapeutic Challenges and Advances}}\\[4pt]
    \textbf{Author Name$^{1}$, Co-author Name$^{2}$}\\
    $^{1}$Department, Institution, Email\\
    $^{2}$Department, Institution, Email\\
    \today
\end{center}

\begin{abstract}
NRAS mutations are key drivers in multiple aggressive cancers, notably melanoma. Despite NRAS being a challenging therapeutic target, recent advances in kinase inhibitors that modulate NRAS downstream signaling provide promising strategies for intervention. This review highlights the biology of NRAS mutations, current kinase inhibitor therapies, combination approaches, and future perspectives. [Placeholder]
\end{abstract}

\textbf{Keywords:} NRAS, kinase inhibitors, melanoma, targeted therapy, MAPK pathway, combination treatments [Placeholder]

\section*{Introduction}
NRAS is a member of the RAS family of small GTPases frequently mutated in melanoma, leukemia, and other cancers. Oncogenic NRAS mutations lead to constitutive activation of signaling pathways controlling proliferation and survival[1][3]. Direct targeting of NRAS remains challenging due to its high affinity for GTP and lack of druggable pockets. Kinase inhibitors targeting downstream effectors of NRAS offer viable therapeutic alternatives. [Placeholder]

\section*{NRAS Biology and Mutation Landscape}
- Overview of NRAS structure, mutation hotspots (e.g., Q61)[1].
- Effects of NRAS mutations on MAPK and PI3K-AKT pathways.
- NRAS-driven tumor phenotypes and clinical implications[2].
[Placeholder]

\section*{Kinase Inhibitors Targeting NRAS-Driven Signaling}
- MEK inhibitors (e.g., trametinib) in NRAS-mutant melanoma; efficacy and resistance patterns[3].
- Other downstream kinase targets: ERK, CDK4/6, FGFR2, and their inhibitors[2][3].
- STK19 kinase as a novel upstream regulator and its inhibitors (e.g., chelidonine)[1].
[Placeholder]

\section*{Combination Therapeutic Strategies}
- Rationale for combining MEK inhibitors with CDK4/6, pan-RAF, or PI3K inhibitors to overcome resistance[3].
- Antisense oligonucleotide approaches to knock down NRAS mRNA in synergy with kinase inhibition[2].
- Kinase activity profiling for personalized combination therapies[2].
[Placeholder]

\section*{Clinical Trials and Emerging Therapies}
- Summary of key clinical trials of kinase inhibitors in NRAS-mutant cancers[3].
- Challenges in therapeutic resistance and toxicity management.
- Prospects for novel agents, immunotherapy combinations, and precision medicine[4][5].
[Placeholder]

\section*{Conclusions and Future Directions}
- Importance of understanding NRAS biology to inform treatment.
- Potential of kinase inhibitor combinations to bypass limitations of single-agent therapies.
- Need for biomarker-driven clinical trials, and exploration of direct NRAS-targeting methods[1][6].
[Placeholder]

% Table placeholder
\begin{table}[h]
\centering
\begin{tabular}{lll}
\hline
Inhibitor & Target & Clinical Status \\
\hline
Trametinib & MEK1/2 & Approved, used in NRAS-mutant melanoma[3] \\
Ribociclib & CDK4/6 & Clinical trials for combination therapies[3] \\
Chelidonine & STK19 kinase & Preclinical, experimental inhibitor[1] \\
[Placeholder] & [Placeholder] & [Placeholder] \\
\hline
\end{tabular}
\caption{Representative kinase inhibitors targeting NRAS signaling pathways. [Placeholder]}
\end{table}

\bibliographystyle{plain}
\bibliography{references}
% Add references as needed.

\end{document}
